\clearpage\begingroup
\HistoricalHead{188}{L}{PHYSIK.\quad NACHLASS.}
\begin{center}9.\end{center}
\noindent\hspace*{\parindent}Die Rechnungen des vorhergehenden Artikels sind hauptsächlich nur als Beispiele zur Erläuterung der Methode zu betrachten. Sie zeigen, dass der
\clearpage\SourceLeaf{194}\HistoricalHead{189}{R}{WECHSELWIRKUNGEN ZWISCHEN GALVANISCHEN STRÖMEN UND MAGNETISMUS.}
\noindent Multiplicator auf die einzelnen Punkte innerhalb desselben ziemlich ungleich wirkt, und es erhellet daraus, dass eine genaue Berechnung der ganzen auf die Nadel wirkenden Kraft die Kenntniss der Vertheilung des Magnetismus in der Nadel voraussetzen würde. Ist $dm$ ein Element dieses Magnetismus, $x$ sein Abstand von einer durch den Aufhängungspunkt der Nadel senkrecht gegen deren magnetische Axe gelegten Vertical-Ebene (welche wir auf der Nordseite als positiv betrachten); $v$ die Multiplicatorkraft in jenem Punkte in der Richtung von Ost nach West senkrecht gegen den magnetischen Meridian, wie sie in den Beispielen des vorigen Art[ikels] berechnet ist; $g$ die Intensität des galvanischen Stromes, welcher den Multiplicator durchläuft, so wird
\[
g\int vxdm
\]
\noindent das Moment der Drehung um eine durch den Aufhängungspunkt gehende Verticallinie. Wäre $v$ in dem ganzen Raume, welchen die Nadel einnimmt, constant, so würde jenes Integral $gv\int xdm$ sein, wo offenbar $\int xdm$ das Moment des Magnetismus der Nadel ausdrückt, welches wir mit $M$ bezeichnen wollen. Setzen wir den horizontalen Theil der Erdmagnetischen Kraft $=T$, so wird die aus dieser und der Elektromagnetischen Kraft des Stromes im Multiplicator zusammengesetzte $=\surd(TT+ggvv)$, und $\frac{gv}{T}$ die Tangente ihrer Neigung gegen den magnetischen Meridian, und in dieser Richtung wird die Nadel im Gleichgewicht sein, insofern $v$ auch innerhalb des Raumes, welchen die Nadel in dieser Lage einnimmt, denselben Werth hat. Alle diese Voraussetzungen bleiben hinreichend genau gültig, wenn sich in der Mitte des Multiplicators eine sehr kleine Nadel befindet. Allein dieses ist nicht mehr der Fall bei der Nadel, wovon hier die Rede ist. Setzt man für diese
\[
\frac{\int vxdm}{\int xdm}\quad\text{oder}\quad\frac{\int vxdm}{M}=v^0,
\]
\noindent so wird $v^0$ als eine Art von Mittelwerth von $v$ angesehen werden können, der von der Vertheilung des Magnetismus in der Nadel abhängig ist, und ohne vollständige Kenntniss von dieser nicht a priori berechnet werden kann: wir werden aber weiter unten eine Methode entwickeln, durch die er aus Combination von Versuchen für unsere Nadel $=8{,}7196$ gefunden ist.

Das Drehungsmoment der Kraft des galvanischen Stromes zur Drehung der Nadel um die verticale Aufhängungsaxe, welches, wenn die Nadel sich im
\clearpage\SourceLeaf{195}\HistoricalHead{190}{L}{PHYSIK.\quad NACHLASS.}
\noindent magnetischen Meridian befindet, $=gv^0M$ ist, wird nun allerdings bei anderer Stellung der Nadel einen anderen Werth haben, dessen Bestimmung eine rein mathematische Aufgabe wird, sobald die Vertheilung des Magnetismus in der Nadel bekannt ist. In Ermangelung dieser Kenntniss können wir jedoch wenigstens übersehen, dass bei einer mässigen Neigung der magnetischen Axe der Nadel gegen den Meridian $=\varphi$, jener Werth durch eine Reihe ausgedrückt werden kann, welche die Form hat
\[
gM(v^0+v'\varphi\varphi+v''\varphi^4+\text{etc.}),
\]
\noindent wenn man, was ohne Zweifel erlaubt ist, den Magnetismus in der Nadel in Beziehung auf ihre Breite symmetrisch vertheilt voraussetzt. Das Drehungsmoment der Nadel durch den Erdmagnetismus wird hingegen, $\varphi$ von der Rechten nach der Linken gezählt,
\[
=-MT\sin\varphi=-MT(\varphi-\tfrac16\varphi^3+\text{etc.})
\]
\noindent sein. Für die Gleichgewichtslage haben wir daher
\[
0=gv^0-T\varphi+gv'\varphi\varphi-\tfrac16T\varphi^3+\text{etc.}
\]
\noindent Es wird mithin der Werth von $\varphi$ für den Ruhestand der Nadel durch die Reihe
\[
\varphi=\frac{gv^0}{T}-\frac{g^3v^0v^0v'}{T^3}\quad\text{etc.}
\]
\noindent ausgedrückt werden, wo man sich jedoch, in sofern es sich nur um solche electromagnetische Kräfte handelt, die nur eine mässige Abweichung hervorbringen können, auf das erste Glied beschränken, oder
\[
\varphi=\frac{gv^0}{T},
\]
\noindent oder, was eben so genau ist,
\[
\operatorname{tang}\varphi=\frac{gv^0}{T}
\]
\noindent setzen kann. Es würde übrigens, wenn man es der Mühe werth hielte, nicht an Mitteln fehlen, auch folgende Coefficienten wenigstens näherungsweise durch Combination solcher Beobachtung auszumitteln, wo man die Nadel unter dem Einfluss des galvanischen Stromes grössere und kleinere Schwingungen machen lässt und das Mittel der Ausweichungen vergleicht. Indessen [er]geben die Versuche, dass man innerhalb des Messungsvermögens unserer Apparate, die auf
\clearpage\SourceLeaf{196}\HistoricalHead{191}{R}{WECHSELWIRKUNGEN ZWISCHEN GALVANISCHEN STRÖMEN UND MAGNETISMUS.}
\noindent Werthe des Winkels $\varphi$ von wenig Geraden beschränkt sind, die folgenden Glieder füglich bei Seite setzen kann.
\clearpage\endgroup
